\documentclass[11pt,letterpaper]{article}
\usepackage[utf8]{inputenc}
\usepackage[T1]{fontenc}
\usepackage[spanish,es-nodecimaldot,es-tabla]{babel}
\usepackage{mathptmx}
\usepackage{amsmath,amssymb}
\usepackage{graphicx}
\usepackage{booktabs,array}
\usepackage{tabularx}
\usepackage{float}
\usepackage{enumitem}
\usepackage[hyphens]{url}
\usepackage[hidelinks]{hyperref}
\usepackage{geometry}
\usepackage{xcolor}
\usepackage[most]{tcolorbox}
\geometry{letterpaper,left=2.4cm,right=2.4cm,top=2.3cm,bottom=2.3cm}
\linespread{1.1}
\setlength{\parindent}{0pt}
\setlength{\parskip}{0.5em}
\setlist[itemize]{leftmargin=1.4em,topsep=1pt,itemsep=1pt}
\setlist[enumerate]{leftmargin=1.7em,topsep=1pt,itemsep=1pt}
\definecolor{iiBlue}{HTML}{1F4E79}
\definecolor{iiOrange}{HTML}{C55A11}
\newtcolorbox{propbox}[1]{colback=iiBlue!6,colframe=iiBlue,coltitle=white,fonttitle=\bfseries,title=#1,arc=2pt,boxrule=0.9pt,left=6pt,right=6pt,top=3pt,bottom=3pt}
\newtcolorbox{alertbox}[1]{colback=iiOrange!7,colframe=iiOrange,coltitle=white,fonttitle=\bfseries,title=#1,arc=2pt,boxrule=0.9pt,left=6pt,right=6pt,top=3pt,bottom=3pt}
\newcommand{\hrulethin}{\noindent\rule{\textwidth}{0.6pt}}
\begin{document}\pagestyle{plain}
\begin{center}
    \textbf{\large Universidad de Santiago de Chile}\\[0.1cm]
    Departamento de Ingeniería Industrial
\end{center}
\vspace{0.15cm}\hrulethin\vspace{0.35cm}
\begin{center}
    {\scshape Estadística Aplicada \quad$\cdot$\quad Caso 1}\\[0.35cm]
    {\LARGE \textbf{Guía del informe (respaldo)}}\\[0.15cm]
    {\large Proyecto Integrador: Regresión}\\[0.2cm]
    Predicción de una respuesta continua
\end{center}
\vspace{0.2cm}\hrulethin\vspace{0.4cm}
El informe es un \textbf{respaldo} del producto evaluado (presentación y script). Máximo \textbf{5 páginas}, sin contar portada ni referencias. Cada tabla o figura debe sustentar una decisión; sin volcados de consola. El código completo vive en el script reproducible.

\section*{Estructura sugerida (6 secciones breves)}
\begin{enumerate}
\item \textbf{Problema, datos y disponibilidad temporal.} Decisión, unidad de análisis, respuesta, población; fuente y licencia; faltantes y su mecanismo; qué información existe al instante de predicción.
\item \textbf{Relaciones e inferencia.} Al menos una relación estimada con intervalo; significancia frente a relevancia; supuestos.
\item \textbf{Validación y preparación.} Partición coherente, validación cruzada, semillas y transformaciones dentro de folds (sin leakage).
\item \textbf{Modelos.} Baseline, lineal, ridge/lasso/elastic-net, splines/GAM y ensembles; hiperparámetros y su selección.
\item \textbf{Comparación, test y subgrupos.} RMSE/MAE/$R^2$ frente al baseline con incertidumbre; test único; error por subgrupo e intervalos de predicción.
\item \textbf{Interpretación, límites y recomendación.} Coeficientes, importancia y PDP/ICE (sin causalidad); extrapolación, drift y recomendación accionable.
\end{enumerate}

\section*{Tabla mínima de validación}
\begin{center}\small
\begin{tabular}{lllllll}
\toprule
Modelo & Familia & RMSE & MAE & $R^2$ & RMSE subgrupo & vs baseline \\
\midrule
Baseline & --- & & & & & \\
OLS & Lineal & & & & & \\
Ridge / Lasso / EN & Regularizado & & & & & \\
Splines / GAM & Flexible & & & & & \\
Random Forest & Ensemble & & & & & \\
Boosting & Ensemble & & & & & \\
Stacking / averaging & Ensemble & & & & & \\
\bottomrule
\end{tabular}\end{center}

El sistema final se reporta aparte (validación vs test), con mejora frente al baseline. Incluya un anexo con el diccionario de variables.

\section*{Preguntas que el informe debe poder responder}
\begin{itemize}
\item ¿Qué decisión apoya la predicción y qué información existe al instante de decidir?
\item ¿Por qué la partición y la validación representan el uso futuro (sin leakage)?
\item ¿El sistema supera el baseline con una magnitud útil y estable?
\item ¿Qué límites y qué monitoreo condicionan su uso?
\end{itemize}\end{document}
